\section*{الفصل \ref{ch:g12:buffers-predominance} --- المحاليل المنظِّمة ومخططات الغلبة}

\begin{solution}{exo:g12:buffers-predominance:1}
\omterm{def:g9:acids-bases-ph:ph}{pH} 2: \ce{CH3COOH}. \omterm{def:g9:acids-bases-ph:ph}{pH} 4.76: الشكلان بكميتين متساويتين. \omterm{def:g9:acids-bases-ph:ph}{pH} 7: \ce{CH3COO-}.
\end{solution}

\begin{solution}{exo:g12:buffers-predominance:2}
\omterm{def:g9:acids-bases-ph:ph}{pH} 4: $1/(1 + 10^{0.76}) = 0.15$. \omterm{def:g9:acids-bases-ph:ph}{pH} 6: $1/(1 + 10^{-1.24}) = 0.95$.
\end{solution}

\begin{solution}{exo:g12:buffers-predominance:3}
أصفر؛ أخضر (داخل \omterm{def:g12:buffers-predominance:indicator}{مجال تغيّر لونه})؛ أزرق.
\end{solution}

\begin{solution}{exo:g12:buffers-predominance:4}
عند \omterm{def:g9:acids-bases-ph:ph}{pH} 7.4، أي تحت 9.26: \ce{NH4+}. وعند \omterm{def:g9:acids-bases-ph:ph}{pH} 11: \ce{NH3}.
\end{solution}

\begin{solution}{exo:g12:buffers-predominance:5}
\omterm{def:g12:buffers-predominance:indicator}{الكاشف الملوّن} نفسه \omterm{def:g12:ka-and-pka:bronsted}{مزدوجة حمض/قاعدة}: فلو استُعمل بكميات كبيرة
لغيّر قيمة \omterm{def:g9:acids-bases-ph:ph}{pH} التي يُراد أن يبيّنها.
\end{solution}

\begin{solution}{exo:g12:buffers-predominance:6}
$10^{\mathrm{pH} - 4.76}$: 0.1; 1; 10.
\end{solution}

\begin{solution}{exo:g12:buffers-predominance:7}
الفينول فثالين (8.0--10.0) أو أزرق الثيمول (8.0--9.1): فمجالاهما
يحويان 8.7. أما أحمر الميثيل فيتغير لونه بين 4.2 و6.3، بعيدًا عن 8.7: وكان
سيتغير قبل ذلك بكثير.
\end{solution}

\begin{solution}{exo:g12:buffers-predominance:8}
$\mathrm{pH} = 4.76 + \log(0.10/0.20) = 4.76 - 0.30 = 4.46$.
\end{solution}

\begin{solution}{exo:g12:buffers-predominance:9}
\omterm{def:g9:acids-bases-ph:ph}{pH} 1: \omterm{def:g9:ions:ion}{الكاتيون} \ce{H3N+-CH2-COOH}، شحنته $+1$. \omterm{def:g9:acids-bases-ph:ph}{pH} 6: \omterm{def:g9:ions:ion}{الأيون} المزدوج،
شحنته الكلية 0. \omterm{def:g9:acids-bases-ph:ph}{pH} 12: \omterm{def:g9:ions:ion}{الأنيون} \ce{H2N-CH2-COO-}، شحنته $-1$.
\end{solution}

\begin{solution}{exo:g12:buffers-predominance:10}
الماء: من 7.0 إلى $-\log(\num{1.0e-4}/0.101) = 3.0$، أي انخفاض بأربع وحدات.
\omterm{def:g12:buffers-predominance:buffer}{المحلول المنظِّم}: من 4.76 إلى $4.76 + \log(9.9/10.1) = 4.75$، أي انخفاض بمقدار 0.01.
\end{solution}

\begin{solution}{exo:g12:buffers-predominance:11}
$\ce{NH4+}/\ce{NH3}$ ($pK_a = 9.26$). والنسبة
$[\ce{NH3}]/[\ce{NH4+}] = 10^{9.5 - 9.26} = 1.7$.
\end{solution}

\begin{solution}{exo:g12:buffers-predominance:12}
تنخفض قيمة \omterm{def:g9:acids-bases-ph:ph}{pH} بوحدة واحدة حين تبلغ النسبة $1/10$:
$(10 - n)/(10 + n) = 0.1$، إذن $n = \qty{8.2}{mmol}$. \omterm{def:g12:buffers-predominance:buffer}{والمحلول المنظِّم} الذي يحوي مزيدًا
من كل شكل يستطيع امتصاص مزيد من الحمض أو القاعدة للتغير نفسه في
النسبة: فسعته أكبر.
\end{solution}

\begin{solution}{exo:g12:buffers-predominance:13}
النسبة $10^{5.0 - 4.76} = 1.74$؛ وكمية الحمض \qty{10}{mmol}، إذن
\qty{17.4}{mmol} من الإيثانوات، أي \qty{174}{mL} من \omterm{def:g2:mixing-and-dissolving:dissolve}{المحلول}.
\end{solution}

\begin{solution}{exo:g12:buffers-predominance:14}
له مزدوجتا حمض/قاعدة، بقيمتين للثابت $pK_a$، ومن ثم تغيّران في
اللون: أحمر عند \omterm{def:g9:acids-bases-ph:ph}{pH} 1، وأصفر عند \omterm{def:g9:acids-bases-ph:ph}{pH} 5، وأزرق عند \omterm{def:g9:acids-bases-ph:ph}{pH} 10.
\end{solution}

\begin{solution}{exo:g12:buffers-predominance:15}
\omterm{def:g10:concentration-and-dilution:dilution}{التخفيف} يقسم التركيزين كليهما على 10: فلا تتغير النسبة، ومن ثم قيمة \omterm{def:g9:acids-bases-ph:ph}{pH}
(4.76). أما \omterm{def:g12:ka-and-pka:strong-acid}{الحمض القوي} المخفف عشر مرات فترتفع قيمة \omterm{def:g9:acids-bases-ph:ph}{pH} له بوحدة
واحدة.
\end{solution}

\begin{solution}{pb:g12:buffers-predominance:1}
\textbf{1.} \ce{CO2 + 2H2O <=> HCO3- + H3O+}.

\textbf{2.} ثنائي أكسيد الكربون \omterm{def:g6:solutions-and-solubility:solute}{المذاب} (مع الماء) هو الحمض،
والهيدروجينوكربونات القاعدة.

\textbf{3.} قيمة الجدول عند \qty{25}{\celsius} في الماء النقي؛ أما
الدم فعند \qty{37}{\celsius} ومليء \omterm{def:g9:ions:ion}{بأيونات} أخرى، تغيّر
الثابت (وعلماء وظائف الأعضاء يحسبون كل ثنائي أكسيد الكربون \omterm{def:g6:solutions-and-solubility:solute}{المذاب}).

\textbf{4.} عند \qty{24}{mmol/L}، $0.024 \times 5.0 = \qty{0.12}{mol}$.

\textbf{5.} محور \omterm{def:g9:acids-bases-ph:ph}{pH} مقسوم عند 6.1: \ce{CO2} تحتها، و\ce{HCO3-} فوقها.

\textbf{6.} عند 7.4، أي فوق 6.1: الهيدروجينوكربونات.

\textbf{7.} قاعدي قليلًا: فقيمة \omterm{def:g9:acids-bases-ph:ph}{pH} فيه أعلى من 7.

\textbf{8.} $10^{7.4 - 6.1} = 10^{1.3} = 20$.

\textbf{9.} $24 / 20 = \qty{1.2}{mmol/L}$.

\textbf{10.} $10^{1.28} = 19$ و$10^{1.32} = 21$.

\textbf{11.} $20/21 = \qty{95}{\percent}$.

\textbf{12.} $10^{7.6 - 6.1} = 10^{1.5} = 32$.

\textbf{13.} \omterm{def:g9:ions:ion}{أيون} الهيدروجينوكربونات:
\ce{HCO3- + H3O+ -> CO2 + 2H2O}.

\textbf{14.} $10^{7.1 - 6.1} = 10$.

\textbf{15.} من $24/20 = 1.2$ إلى $24/10 = \qty{2.4}{mmol/L}$: أي يتضاعف.
فتتنفس الرئتان أسرع وأعمق، لطرد ثنائي أكسيد الكربون
الزائد.

\textbf{16.} يستطيع \omterm{def:g12:ka-and-pka:strong-acid}{الحمض الضعيف} وقاعدته أن يمتص كل منهما ما يضاف،
ولا تتوقف قيمة \omterm{def:g9:acids-bases-ph:ph}{pH} إلا على النسبة بينهما؛ أما \omterm{def:g12:ka-and-pka:strong-acid}{الحمض القوي} فلا يفعل إلا أن يدفع
قيمة \omterm{def:g9:acids-bases-ph:ph}{pH} في اتجاه واحد.

\textbf{17.} لا تتغير إطلاقًا: فالنسبة، ومن ثم قيمة \omterm{def:g9:acids-bases-ph:ph}{pH}، تبقيان على حالهما.

\textbf{18.} نحو 20.
\end{solution}
